% !TeX program = xelatex
% Звіт з комп'ютерного практикуму №1, Київ, 2026.
% Компіляція: xelatex -interaction=nonstopmode -halt-on-error report.tex (двічі)
% Скріншоти зберігаються у page-screenshots/Desktop та page-screenshots/Mobile.
\documentclass[12pt,a4paper]{article}
\usepackage{iftex}
\ifPDFTeX
  \usepackage[T2A]{fontenc}
  \usepackage[utf8]{inputenc}
\else
  \usepackage{fontspec}
  \setmainfont{Liberation Serif}
  \setsansfont{Noto Sans}
  \setmonofont{DejaVu Sans Mono}[Scale=0.84]
\fi
\usepackage[ukrainian]{babel}
\usepackage[a4paper,left=23mm,right=22mm,top=23mm,bottom=24mm]{geometry}
\usepackage{microtype}
\usepackage{setspace}
\setstretch{1.17}
\setlength{\emergencystretch}{3em}
\usepackage{graphicx}
\usepackage{array,longtable,tabularx,booktabs}
\usepackage{enumitem}
\usepackage{xcolor}
\usepackage{tcolorbox}
\usepackage{tikz}
\usepackage{float}
\usepackage{caption}
\usepackage{fancyhdr}
\usepackage{titlesec}
\usepackage{hyperref}
\usepackage{url}
\usepackage{lastpage}
\usepackage{needspace}
\usepackage{ragged2e}

\newcommand{\AuthorName}{Суворов Дмитро}
\newcommand{\StudentGroup}{ЦІ-31}
\newcommand{\ProjectRepo}{\url{https://github.com/Dima018/Gym_Website}}
\newcommand{\ProjectBoard}{\url{https://github.com/users/Dima018/projects/13}}

\definecolor{forest}{HTML}{233C30}
\definecolor{sage}{HTML}{5B7957}
\definecolor{mist}{HTML}{F0F4EE}
\definecolor{rulegray}{HTML}{D9E1DA}
\hypersetup{hypertexnames=false,colorlinks=true,linkcolor=forest,urlcolor=sage,pdftitle={Звіт КП №1 — сайт спортзалу Ритм},pdfauthor={Суворов Дмитро}}
\titleformat{\section}{\Large\sffamily\bfseries\color{forest}}{\thesection.}{0.6em}{}
\titleformat{\subsection}{\large\sffamily\bfseries\color{forest}}{\thesubsection.}{0.6em}{}
\titlespacing*{\section}{0pt}{15pt}{9pt}
\titlespacing*{\subsection}{0pt}{12pt}{6pt}
\captionsetup{hypcap=false,font=small,labelfont={bf,color=forest},justification=centering}
\setlist[itemize]{leftmargin=1.45em,itemsep=2pt,topsep=4pt}
\setlist[enumerate]{leftmargin=1.55em,itemsep=2pt,topsep=4pt}
\renewcommand{\arraystretch}{1.20}
\pagestyle{fancy}
\fancyhf{}
\fancyhead[L]{\small\sffamily\color{sage} КП №1 · WEB-орієнтовані системи}
\fancyhead[R]{\small\sffamily\color{sage} Ритм}
\fancyfoot[L]{\footnotesize\color{sage} \AuthorName\ · 2026}
\fancyfoot[R]{\footnotesize\color{sage} \thepage\,/\,\pageref{LastPage}}
\renewcommand{\headrulewidth}{0.4pt}
\renewcommand{\headrule}{\hbox to\headwidth{\color{rulegray}\leaders\hrule height \headrulewidth\hfill}}
\setlength{\headheight}{26pt}
\newtcolorbox{infobox}{colback=mist,colframe=rulegray,boxrule=0.5pt,arc=2mm,left=3mm,right=3mm,top=2mm,bottom=2mm}
\newcommand{\file}[1]{\texttt{#1}}
\newcommand{\evidencehead}[2]{%
  \clearpage\subsection*{#1 --- #2}\phantomsection\addcontentsline{toc}{subsection}{#1 --- #2}
}
\newcommand{\deskimage}[2]{%
 \begin{center}\includegraphics[width=\linewidth,height=10.1cm,keepaspectratio]{page-screenshots/#1}\par\vspace{2pt}{\footnotesize\color{sage} #2}\end{center}
}
\newcommand{\mobileimage}[2]{%
 \begin{minipage}[t]{0.480\linewidth}\centering
 \fcolorbox{rulegray}{white}{\includegraphics[width=0.96\linewidth,height=18.4cm,keepaspectratio]{page-screenshots/#1}}\par\smallskip{\footnotesize\color{sage}#2}\end{minipage}
}
\begin{document}
\begin{titlepage}
\thispagestyle{empty}
\begin{center}
{\small\sffamily\bfseries МІНІСТЕРСТВО ОСВІТИ І НАУКИ УКРАЇНИ}\par\vspace{5mm}
{\small\sffamily\bfseries НАЦІОНАЛЬНИЙ ТЕХНІЧНИЙ УНІВЕРСИТЕТ УКРАЇНИ}\par
{\small\sffamily\bfseries «КИЇВСЬКИЙ ПОЛІТЕХНІЧНИЙ ІНСТИТУТ ІМЕНІ ІГОРЯ СІКОРСЬКОГО»}\par
\vspace{25mm}
{\color{sage}\rule{38mm}{1.8pt}}\par\vspace{8mm}
{\LARGE\sffamily\bfseries\color{forest} ЗВІТ}\par\vspace{5mm}
{\Large\sffamily Комп'ютерний практикум №1}\par\vspace{3mm}
{\large «Розробка простого веб-сайту зі стилізацією. HTML. CSS.»}\par\vspace{11mm}
{\large з дисципліни}\par\vspace{1mm}
{\large\bfseries «Технології розробки WEB-орієнтованих систем»}\par\vspace{12mm}
\begin{tcolorbox}[width=0.87\textwidth,colback=mist,colframe=rulegray,arc=2mm,boxrule=0.7pt]
\centering{\large\sffamily\bfseries\color{forest} Тема №69: «Спортзал»}\par\smallskip
{\normalsize Навчальний веб-сайт спортзалу «Ритм»}
\end{tcolorbox}
\vfill
\begin{flushright}
\begin{tabular}{@{}rl@{}}
Виконавець: & \textbf{\AuthorName}\\
Група: & \StudentGroup\\
Викладач: & С. В. Суліма\\
\end{tabular}
\end{flushright}
\vfill
{\large Київ --- 2026}
\end{center}
\end{titlepage}

\tableofcontents
\clearpage

\section{Тема, мета та завдання}
\textbf{Тема:} багатосторінковий веб-сайт спортивного залу «Ритм» (тема №69).

\textbf{Мета:} набути практичних навичок розроблення семантичної HTML5-розмітки, оформлення сторінок засобами CSS3, створення адаптивного інтерфейсу та використання вбудованої перевірки HTML-форм. Робота відповідає першому етапу наскрізного веб-проєкту, без серверної частини та бази даних.

\textbf{Основні завдання:}
\begin{enumerate}
\item Спроєктувати структуру сайту та ескізи основних сторінок.
\item Реалізувати п'ять взаємопов'язаних HTML-сторінок зі спільною навігацією.
\item Стилізувати інтерфейс за допомогою CSS, Flexbox та Grid.
\item Забезпечити адаптивне відображення на настільних і мобільних екранах.
\item Створити навчальну форму запису з браузерною валідацією.
\item Підготувати скриншоти й матеріали для перевірки виконаної роботи.
\end{enumerate}

\section{Авторство та організація роботи}
\noindent\begin{tabularx}{\textwidth}{>{\RaggedRight\arraybackslash\bfseries}p{38mm}X}
\toprule
Виконавець & \AuthorName\\
Профіль GitHub & \href{https://github.com/Dima018}{github.com/Dima018}\\
Група & \StudentGroup\\
Склад бригади & Суворов Дмитро\\
Розподіл ролей & \AuthorName: розроблення сайту, перевірка та підготовка звіту\\
\bottomrule
\end{tabularx}

\section{Структура веб-сайту}
Сайт містить п'ять сторінок, які пов'язані спільним меню, а також посиланнями між картками тренувань, сторінкою детального опису та формою запису.

\begin{center}
\begin{tikzpicture}[node distance=9mm and 10mm, every node/.style={font=\small\sffamily}, card/.style={draw=rulegray,fill=mist,rounded corners=2mm,align=center,minimum height=10mm,inner sep=3mm}]
\node[card] (home) at (0,0) {Головна\\\texttt{index.html}};
\node[card] (train) at (-4,-2) {Каталог тренувань\\\texttt{trainings.html}};
\node[card] (about) at (0,-2) {Про спортзал\\\texttt{about-gym.html}};
\node[card] (sign) at (4,-2) {Форма запису\\\texttt{training-signup.html}};
\node[card] (strength) at (-4,-4) {Силове тренування\\\texttt{strength-training.html}};
\draw[->,thick,sage] (home) -- (train);
\draw[->,thick,sage] (home) -- (about);
\draw[->,thick,sage] (home) -- (sign);
\draw[->,thick,sage] (train) -- (strength);
\draw[->,thick,sage] (train) -- (sign);
\draw[->,thick,sage] (strength) -- (sign);
\end{tikzpicture}
\captionof{figure}{Карта сторінок сайту «Ритм» (основні переходи)}
\end{center}

\noindent\begin{tabularx}{\textwidth}{>{\ttfamily\small}p{49mm}X}
\toprule
\multicolumn{1}{l}{\textbf{Файл}} & \textbf{Призначення}\\
\midrule
index.html & Головна сторінка, презентація залу та переваг\\
trainings.html & Каталог шести напрямів тренувань із цінами\\
strength-training.html & Детальний опис силового заняття, програма й розклад\\
about-gym.html & Інформація про зал, контакти й години роботи\\
training-signup.html & Форма запису з HTML5-перевіркою полів\\
styles/main.css & Спільний дизайн та правила адаптивності\\
assets/ & Ілюстрації та графічні матеріали\\
assets/logo.png & Логотип спортзалу «Ритм»\\
\bottomrule
\end{tabularx}

\section{Реалізація веб-сайту}
\subsection{HTML5 та семантична структура}
Сторінки мають заголовок документа, метадані, область навігації, основний вміст і підвал. Для цього використано елементи \file{header}, \file{nav}, \file{main}, \file{section}, \file{article} та \file{footer}. Інформаційні блоки відокремлено логічно, сторінки мають виразний заголовок першого рівня. Для доступності передбачено альтернативні підписи \file{alt} для зображень, підписи \file{label} для полів форми та посилання для переходу до основного вмісту.

\subsection{Оформлення та компонування}
Інтерфейс має світлий фон, темно-зелену кольорову гаму, контрастні кнопки та картки. Основні технічні рішення наведено в таблиці~\ref{tab:css}.
\begin{table}[H]
\centering
\caption{Ключові прийоми стилізації}\label{tab:css}
\noindent\begin{tabularx}{\textwidth}{>{\RaggedRight\arraybackslash\bfseries}p{40mm}X}
\toprule
Прийом & Застосування\\
\midrule
CSS Grid & Три картки каталогу в рядку, двоколонкові інформаційні секції\\
Flexbox & Навігація, кнопкові блоки та взаємне вирівнювання елементів\\
Media queries & Зміна компонування на ширинах 900 і 600 пікселів\\
Відносні розміри & Контейнер із максимальною шириною та гнучкі зображення\\
Графіка та шрифти & Оптимізовані локальні зображення та системний стек шрифтів\\
Псевдокласи & Стани наведення, клавіатурного фокусу й цільового блока\\
\bottomrule
\end{tabularx}
\end{table}

\subsection{Адаптивний дизайн}
За великої ширини в каталозі відображено три картки в ряд; на проміжних екранах --- дві, а на вузьких --- одну. Двоколонкові блоки головної сторінки, опису тренування, контактів та форми на мобільних пристроях переходять до вертикального розташування. Меню переноситься на наступний рядок замість горизонтального переповнення. На зображеннях у розділі~\ref{sec:screenshots} показано обидва режими.

\subsection{Навчальна форма запису}
Сторінка \file{training-signup.html} включає поля імені, електронної пошти, телефону, напряму занять, довільного коментаря та підтвердження ознайомлення з навчальним характером форми. Для перевірки використано стандартні можливості HTML5: атрибути \file{required}, \file{minlength}, \file{maxlength}, \file{pattern} та \file{type="email"}. Поле телефону приймає формат \file{+380} і дев'ять цифр. Після успішної локальної перевірки відкривається блок результату через CSS \file{:target}.

Форма \textbf{не передає дані на сервер} й не створює реального запису на тренування, що відповідає статичному етапу практикуму №1.

\section{Перевірка результатів}
Для фінальної версії виконано локальну перевірку верстки в Chromium на ширинах 320, 390, 768 та 1440 пікселів. Сторінки відкривалися через локальний HTTP-сервер; перевірено завантаження зображень і відносні шляхи до ресурсів. Для 20 перевірених варіантів не виявлено горизонтального переповнення, а всі ілюстрації завантажилися. Кожна сторінка має один заголовок \file{h1}. Скріншоти фінальної версії демонструють компонування елементів на настільному та мобільному екранах.

\noindent\begin{tabularx}{\textwidth}{>{\RaggedRight\arraybackslash\bfseries}p{48mm}X}
\toprule
Перевірка & Результат / обмеження\\
\midrule
Відображення сторінок & Представлені desktop та mobile скриншоти для усіх п'яти сторінок\\
Верстка у Chromium & Локальна перевірка Chromium: 5 сторінок, 4 ширини, без горизонтального переповнення\\
HTML5-форма & Наявні обов'язкові поля, перевірка email і шаблону телефону\\
W3C HTML/CSS & Офіційний результат валідаторів не наданий; не позначається як виконаний\\
Публікація коду & Посилання на GitHub-репозиторій наведено у розділі «Посилання та матеріали»\\
Дошка задач & Посилання на GitHub Projects наведено у розділі «Посилання та матеріали»\\
\bottomrule
\end{tabularx}

\section{Скріншоти веб-сайту}\label{sec:screenshots}
Нижче наведено скриншоти фінальної версії сайту з логотипом, зроблені в Chromium на настільних і мобільних розмірах екрана. Зображення взято з папки \file{page-screenshots}. Кожна сторінка представлена обома варіантами відображення.

\evidencehead{Головна сторінка}{desktop}
\deskimage{Desktop/dHome1.png}{Верхня частина: навігація, вступний блок, ілюстрація}
\deskimage{Desktop/dHome2.png}{Нижня частина: переваги, інформаційний банер, підвал}
\evidencehead{Головна сторінка}{mobile}
\noindent\mobileimage{Mobile/mHome1.png}{Верхня частина}\hfill\mobileimage{Mobile/mHome2.png}{Нижня частина}

\evidencehead{Каталог тренувань}{desktop}
\deskimage{Desktop/dTrainings1.png}{Перші картки каталогу}
\deskimage{Desktop/dTrainings2.png}{Другий ряд карток: бокс, стретчинг та персональне тренування}
\evidencehead{Каталог тренувань}{mobile}
\noindent\mobileimage{Mobile/mTrainings1.png}{Верхня частина каталогу}\hfill\mobileimage{Mobile/mTrainings2.png}{Нижня частина каталогу}

\evidencehead{Силове тренування}{desktop}
\deskimage{Desktop/dStrength1.png}{Заголовок, ілюстрація та інформація про заняття}
\deskimage{Desktop/dStrength2.png}{Програма заняття, необхідні речі та розклад}
\evidencehead{Силове тренування}{mobile}
\noindent\mobileimage{Mobile/mStrength1.png}{Верхня частина}\hfill\mobileimage{Mobile/mStrength2.png}{Нижня частина}

\evidencehead{Сторінка «Про нас»}{desktop}
\deskimage{Desktop/dAbout1.png}{Опис спортивного залу та контактна інформація}
\evidencehead{Сторінка «Про нас»}{mobile}
\noindent\mobileimage{Mobile/mAbout1.png}{Вступ і опис}\hfill\mobileimage{Mobile/mAbout2.png}{Контакти та графік роботи}

\evidencehead{Форма запису}{desktop}
\deskimage{Desktop/dSignup1.png}{Основні поля форми та навчальні пояснення}
\deskimage{Desktop/dSignup2.png}{Нижня частина форми та кнопка перевірки}
\evidencehead{Форма запису}{mobile}
\noindent\mobileimage{Mobile/mSignup1.png}{Поля для контактних даних}\hfill\mobileimage{Mobile/mSignup2.png}{Напрям тренувань, коментар і підтвердження}

\clearpage
\section{Висновки}
У ході комп'ютерного практикуму №1 створено багатосторінковий статичний веб-сайт «Ритм» для презентації спортивного залу. Реалізовано п'ять HTML5-сторінок зі спільним меню, каталогом шести напрямів тренувань, детальною інформацією про силове заняття, контактною сторінкою та навчальною формою запису. Єдине CSS-оформлення забезпечує адаптивне розташування елементів, а вбудовані механізми HTML5 --- базову перевірку введених даних.

Скріншоти підтверджують оформлення настільних і мобільних версій усіх сторінок. Виконано локальну перевірку верстки для чотирьох ширин екрана та валідації навчальної форми. Посилання на репозиторій і дошку задач наведено нижче.

\section{Посилання та матеріали}
\begin{enumerate}
\item \textbf{Git-репозиторій:} \ProjectRepo
\item \textbf{Профіль автора:} \href{https://github.com/Dima018}{github.com/Dima018}.
\item \textbf{Дошка задач:} \ProjectBoard
\item \textbf{Методичні вказівки:} С.~В.~Суліма, «Методичні вказівки до комп'ютерних практикумів з дисципліни ``Технології розробки WEB-орієнтованих систем''», КПІ ім. Ігоря Сікорського, 2026.
\item \textbf{Локальні матеріали:} файл README, HTML-сторінки, CSS, графічні матеріали та скриншоти сторінок.
\end{enumerate}
\end{document}
